\documentclass[a4paper,11pt]{article}

\usepackage[a4paper,margin=2cm]{geometry}
\usepackage[T1]{fontenc}
\usepackage[utf8]{inputenc}
\usepackage[ngerman,english]{babel}
\usepackage{lmodern}
\usepackage{microtype}
\usepackage{xcolor}
\usepackage{parskip}
\usepackage{enumitem}
\usepackage[most]{tcolorbox}
\usepackage{titlesec}
\usepackage[hidelinks]{hyperref}

\definecolor{HeaderBlueDark}{RGB}{70,110,170}
\definecolor{RowBlueLight}{RGB}{235,243,255}
\definecolor{RowBlueDark}{RGB}{220,233,250}
\definecolor{SoftGray}{RGB}{90,90,90}

\setlength{\parindent}{0pt}
\setlist[itemize]{leftmargin=1.4em,itemsep=0.35em,topsep=0.35em}
\linespread{1.06}

\titleformat{\section}[block]
{\Large\bfseries\color{HeaderBlueDark}}
{}{0pt}{}
\titlespacing{\section}{0pt}{1.3em}{0.8em}

\titleformat{\subsection}[block]
{\large\bfseries\color{HeaderBlueDark}}
{}{0pt}{}
\titlespacing{\subsection}{0pt}{1.0em}{0.6em}

\newtcolorbox{exbox}{enhanced,breakable,colback=RowBlueLight,colframe=RowBlueDark,boxrule=0.4pt,arc=1.5mm,left=1.2mm,right=1.2mm,top=1mm,bottom=1mm,title={Beispiele \textnormal{(Examples)}},fonttitle=\bfseries,colbacktitle=RowBlueDark,coltitle=black}

\NewDocumentEnvironment{grammarentry}{mm+b}{%
    \begin{tcolorbox}[enhanced,breakable,colback=white,colframe=HeaderBlueDark,boxrule=0.6pt,arc=1.5mm,left=1.5mm,right=1.5mm,top=1mm,bottom=1mm,colbacktitle=HeaderBlueDark,coltitle=white,fonttitle=\bfseries,title={#1\hfill\normalfont\footnotesize #2}]
    #3
    \end{tcolorbox}
}{}

\newcommand{\GermanText}[1]{\textbf{Deutsch: }#1\par}
\newcommand{\EnglishText}[1]{{\small\color{SoftGray}\textbf{English: }#1\par}}
\newcommand{\ExampleItem}[2]{\item \textbf{#1}\\{\small\color{SoftGray}#2}}

\title{Substantivierte Adjektive}
\date{}

\begin{document}
    
    \maketitle
    \tableofcontents
    \newpage
    
    
    \section{Was sind substantivierte Adjektive?}
        
        \begin{grammarentry}{Definition}{Definition}
            
            \GermanText{
                Ein substantiviertes Adjektiv ist ein Adjektiv, das wie ein Nomen verwendet wird.
                Es wird deshalb gro{\ss}geschrieben und kann Artikel, Pronomen oder andere Begleiter haben.
            }
            
            \EnglishText{
                A substantivized adjective is an adjective that is used as a noun.
                Therefore, it is written with a capital letter and can be used with articles, pronouns, and determiners.
            }
            
            \GermanText{
                Beispiele:
                \begin{itemize}
                    \item gut $\rightarrow$ das Gute
                    \item neu $\rightarrow$ etwas Neues
                    \item alt $\rightarrow$ der Alte
                    \item jung $\rightarrow$ die Junge
                \end{itemize}
            }
            
            \EnglishText{
                Examples:
                \begin{itemize}
                    \item good $\rightarrow$ the good
                    \item new $\rightarrow$ something new
                    \item old $\rightarrow$ the old man
                    \item young $\rightarrow$ the young woman
                \end{itemize}
            }
        
        \end{grammarentry}
    
    
    \section{Warum sagt man ``etwas Abstraktes''?}
        
        \begin{grammarentry}{Etwas + substantiviertes Adjektiv}{Something + substantivized adjective}
            
            \GermanText{
                Nach W\"ortern wie \textit{etwas}, \textit{nichts}, \textit{viel}, \textit{wenig} und \textit{manches}
                steht oft ein substantiviertes Adjektiv.
            }
            
            \EnglishText{
                After words such as \textit{something}, \textit{nothing}, \textit{much},
                \textit{little}, and \textit{some}, German often uses a substantivized adjective.
            }
            
            \GermanText{
                Deshalb sagt man:
                \begin{itemize}
                    \item etwas Neues
                    \item etwas Interessantes
                    \item etwas Sch\"ones
                    \item etwas Abstraktes
                \end{itemize}
            }
            
            \EnglishText{
                Therefore German says:
                \begin{itemize}
                    \item something new
                    \item something interesting
                    \item something beautiful
                    \item something abstract
                \end{itemize}
            }
            
            \GermanText{
                In dem Satz \glqq Leitung ist oft etwas Abstraktes\grqq{}
                ist \glqq Abstraktes\grqq{} kein Genitiv.
                Es ist ein substantiviertes Adjektiv.
            }
            
            \EnglishText{
                In the sentence ``Leitung ist oft etwas Abstraktes'',
                ``Abstraktes'' is not a genitive form.
                It is a substantivized adjective.
            }
        
        \end{grammarentry}
    
    
    \section{Typische Beispiele}
        
        \begin{exbox}
            
            \begin{itemize}
                
                \ExampleItem
                {Ich habe etwas Neues gelernt.}
                {I learned something new.}
                
                \ExampleItem
                {Heute ist nichts Besonderes passiert.}
                {Nothing special happened today.}
                
                \ExampleItem
                {Das Gute an Deutsch ist die klare Grammatik.}
                {The good thing about German is the clear grammar.}
                
                \ExampleItem
                {Sie sucht etwas Interessantes zum Lesen.}
                {She is looking for something interesting to read.}
                
                \ExampleItem
                {Leitung ist oft etwas Abstraktes.}
                {Leadership or management is often something abstract.}
            
            \end{itemize}
        
        \end{exbox}
    
    
    \section{Substantivierte Adjektive mit Artikeln}
        
        \begin{grammarentry}{Mit bestimmtem Artikel}{With definite article}
            
            \GermanText{
                Substantivierte Adjektive k\"onnen auch mit Artikeln verwendet werden.
            }
            
            \EnglishText{
                Substantivized adjectives can also be used with articles.
            }
            
            \begin{exbox}
                
                \begin{itemize}
                    
                    \ExampleItem
                    {Der Alte wohnt nebenan.}
                    {The old man lives next door.}
                    
                    \ExampleItem
                    {Die Junge liest ein Buch.}
                    {The young woman is reading a book.}
                    
                    \ExampleItem
                    {Das Gute bleibt bestehen.}
                    {The good remains.}
                    
                    \ExampleItem
                    {Die Reichen zahlen viele Steuern.}
                    {The rich pay many taxes.}
                    
                    \ExampleItem
                    {Die Armen brauchen Hilfe.}
                    {The poor need help.}
                
                \end{itemize}
            
            \end{exbox}
        
        \end{grammarentry}
    
    
    \section{Wichtiger Hinweis}
        
        \begin{grammarentry}{Merke}{Remember}
            
            \GermanText{
                Wenn ein Adjektiv als Nomen benutzt wird,
                schreibt man es gro{\ss}.
                
                Beispiele:
                \begin{itemize}
                    \item etwas Neues
                    \item nichts Wichtiges
                    \item das Gute
                    \item die Reichen
                    \item der Alte
                \end{itemize}
            }
            
            \EnglishText{
                When an adjective is used as a noun,
                it is written with a capital letter.
                
                Examples:
                \begin{itemize}
                    \item something new
                    \item nothing important
                    \item the good
                    \item the rich
                    \item the old man
                \end{itemize}
            }
        
        \end{grammarentry}

\end{document}
